%~Mouliné par MaN_auto v.0.33.4 2025-04-02 14:33:00
\documentclass[AHL,Unicode,longabstracts,published]{cedram}




\newcommand{\slim}{\operatorname{s-}\hspace{-0.25cm}\lim_{t\to\pm\infty}}

\newcommand{\Schr}{Schr\"odinger }

\renewcommand{\a}{\alpha}
\renewcommand{\b}{\beta}
\newcommand{\g}{\gamma}
\renewcommand{\d}{\delta}
\newcommand{\e}{\varepsilon}
\newcommand{\f}{\varphi}
\newcommand{\ps}{\psi}
\newcommand{\z}{\zeta}

\renewcommand{\l}{\lambda}

\newcommand{\x}{\xi}
\newcommand{\y}{\eta}
\newcommand{\s}{\sigma}


\newcommand{\C}{\mathbb{C}}
\newcommand{\R}{\mathbb{R}}
\newcommand{\Z}{\mathbb{Z}}
\newcommand{\T}{\mathbb{T}}
\newcommand{\bbN}{\mathbb{N}}
\newcommand{\bbA}{\mathbb{A}}

\newcommand{\HS}{\mathcal{H}}
\newcommand{\JS}{\mathcal{J}}
\newcommand{\LS}{\mathcal{L}}
\newcommand{\US}{\mathcal{U}}
\newcommand{\FS}{\mathcal{F}}
\newcommand{\BS}{\mathcal{B}}
\newcommand{\CS}{\mathcal{C}}
\newcommand{\DS}{\mathcal{D}}

\newcommand{\ac}{\mathrm{ac}}
\newcommand{\IK}{\mathrm{IK}}
\newcommand{\pp}{\mathrm{pp}}
\newcommand{\la}{\langle}
\newcommand{\ra}{\rangle}
\newcommand{\jap}[1]{\langle #1 \rangle}

\DeclareMathOperator{\kerr}{Ker}
\DeclareMathOperator{\supp}{supp}
\DeclareMathOperator{\Ran}{Ran}
\DeclareMathOperator{\Ferm}{Ferm}
\DeclareMathOperator{\Idd}{Id}


\newenvironment{alcases}{\left\{\begingroup\arraycolsep=0pt\renewcommand{\arraystretch}{1.2}\begin{array}{r@{\;}l@{\quad}l}}{\end{array}\endgroup\right.}

\newcounter{stepcount}
\newenvironment{step}{\refstepcounter{stepcount}\begin{proof}[Step~\thestepcount]}{\let\qed\relax\end{proof}}

%%%%%%%%%%%%%%%%%%%%%%%%%%%%%%%%%%%%%%%%%%%%%%%%%%%%%%%%%%%%%%%%%%%%%%%%%%%%%%%%%%%%%%%
\graphicspath{{./figures/}}

\newcommand*{\mk}{\mkern -1mu}
\newcommand*{\Mk}{\mkern -2mu}
\newcommand*{\mK}{\mkern 1mu}
\newcommand*{\MK}{\mkern 2mu}

\hypersetup{urlcolor=purple, linkcolor=blue, citecolor=red}
\newcommand*{\relabel}{\renewcommand{\labelenumi}{(\theenumi)}}
\newcommand*{\romanenumi}{\renewcommand*{\theenumi}{\roman{enumi}}\relabel}
\newcommand*{\Romanenumi}{\renewcommand*{\theenumi}{\Roman{enumi}}\relabel}
\newcommand*{\alphenumi}{\renewcommand*{\theenumi}{\alph{enumi}}\relabel}
\newcommand*{\Alphenumi}{\renewcommand*{\theenumi}{\Alph{enumi}}\relabel}
%\let\oldtilde\tilde
%\renewcommand*{\tilde}[1]{\mathchoice{\widetilde{#1}}{\widetilde{#1}}{\oldtilde{#1}}{\oldtilde{#1}}}
\let\oldhat\hat
\renewcommand*{\hat}[1]{\mathchoice{\widehat{#1}}{\widehat{#1}}{\oldhat{#1}}{\oldhat{#1}}}
\let\oldforall\forall
\renewcommand*{\forall}{\mathrel{\oldforall}}
%%%%%%%%%%%%%%%%%%%%%%%%%%%%%%%%%%%%%%%%%%%%%%%%%%%%%%%%%%%%%%%%%%%%%%%%%%%%%%%%%%%%%%%

\title[Isozaki--Kitada modifiers on general lattices]{Construction of Isozaki--Kitada modifiers for discrete Schr\"odinger operators on general lattices}
\alttitle{Construction de modificateurs d'Isozaki--Kitada pour les opérateurs de Schr\"odinger discrets sur des réseaux généraux}


\subjclass{47A40, 47B39, 81U05}
\keywords{long-range scattering theory, discrete Schr\"odinger operators, modified wave operators, time-independent modifiers}



\author[\initial{Y.} \lastname{Tadano}]{\firstname{Yukihide} \lastname{Tadano}}
\address{Laboratory of\\
Mathematical Science,\\
Graduate School of Science,\\
University of Hyogo,\\
2167 Shosha, Himeji,\\
Hyogo 671-2201 (Japan)}
\email{tadano@sci.u-hyogo.ac.jp}
\thanks{The author was partially supported by JSPS Grant Numbers 20J00247, 21K20337 and 23K12991. This work is supported by the Research Institute for Mathematical Sciences, an International Joint Usage/Research Center located in Kyoto University. The author would like to thank Shu Nakamura for encouraging me to write this article.}
\CDRGrant{20J00247}
\CDRGrant{21K20337}
\CDRGrant{23K12991}
\begin{abstract}
We consider a scattering theory for difference operators on $\HS=\ell^2(\Z^d; \mathbb{C}^n)$ perturbed with a long-range potential $V:\Z^d\to\R^n$. One of the motivating examples is discrete Schr\"odinger operators on $\Z^d$-periodic graphs. We construct time-independent modifiers, so-called Isozaki--Kitada modifiers, and we prove that the modified wave operators with the above-mentioned Isozaki--Kitada modifiers exist and that they are complete.
\end{abstract}

\begin{altabstract}
Nous considérons une théorie du scattering pour les opérateurs aux différences sur $\HS=\ell^2(\Z^d ; \mathbb{C}^n)$ perturbés par un potentiel à longue portée $V:\Z^d\to\R^n$. Un des exemples motivants est celui des opérateurs de Schr\"odinger discrets sur des graphes $\Z^d$-périodiques. Nous construisons des modificateurs indépendants du temps, appelés modificateurs d'Isozaki--Kitada, et nous prouvons que les opérateurs d'onde modifiés avec les modificateurs d'Isozaki--Kitada mentionnés ci-dessus existent et qu'ils sont complets.
\end{altabstract}

\datereceived{2023-06-06}
\daterevised{2024-08-14}
\dateaccepted{2024-09-30}

\editors{S. Vu Ngoc and N. Anantharaman}
\begin{DefTralics}
\newcommand{\Z}{\mathbb{Z}}
\newcommand{\R}{\mathbb{R}}
\newcommand{\HS}{\mathcal{H}}
\end{DefTralics}
%%%%%%%%%%%%%%%%%%%%%%%%%%%%%%%%%%%%%%%%%%%%%%%%%%%%%%%%%%%%%%%%%%%%%%%%%%%%%%%%%%%%%%%%%%%%%%%%
\dateposted{2025-05-20}
\begin{document}

\maketitle

\section{Introduction}

The aim of the present article is to construct a long-range scattering theory for difference operators on the space of vector-valued functions on $\Z^d$. This problem is motivated by discrete Schr\"odinger operators on an arbitrary non-primitive lattice, e.g., hexagonal lattice, diamond lattice, Kagome lattice and graphite (see~\cite{An-I-M} for more examples). Note that the cases of primitive lattices and the hexagonal lattice are considered in~\cite{T1,T2}, respectively.

Let $\HS=\ell^2(\Z^d; \C^n)$, where $d$ and $n$ are positive integers. For $u \in \HS$, we use the notation
\begin{align*}
u=
\begin{pmatrix}
u_1 \\ u_2 \\ \vdots \\ u_n
\end{pmatrix}
, \quad u_j \in \ell^2\left(\Z^d\right)=\ell^2\left(\Z^d;\C\right).
\end{align*}
We consider a generalized form of discrete \Schr operators on $\HS$:
\begin{align*}
H=H_0+V.
\end{align*}
The unperturbed operator $H_0$ is defined as a convolution operator by $(f_{jk})_{1\,\leq\,j,\,k\,\leq\,n}$, that is,
\begin{align*}
H_0u &= 
\begin{pmatrix}
H_{0,11} & H_{0,12} & \cdots & H_{0,1n} \\
H_{0,21} & H_{0,22} & \cdots & H_{0,2n} \\
\vdots & \vdots & \ddots & \vdots \\
H_{0,n1} & H_{0,n2} & \cdots & H_{0,nn}
\end{pmatrix}
u, \quad u \in \HS, \\
H_{0,jk}u_k(x) &=\sum_{y\,\in\,\Z^d}f_{jk}(x-y)u_k(y), \quad u_k \in \ell^2\left(\Z^d\right).
\end{align*}
Here each $f_{jk}: \Z^d \rightarrow \C$ is a rapidly decreasing function, i.e.,
\begin{align*}
\sup_{x\,\in\,\Z^d} \la x \ra^m \left|f_{jk}(x)\right| <\infty
\end{align*}
for any $m\in\bbN$, where $\la x \ra = (1+|x|^2)^{\frac{1}{2}}$. The perturbation $V$ is a multiplication operator by $V={}^t(V_1,\cdots,V_n):\Z^d \to \R^n$,
\begin{align*}
Vu(x) &= 
\begin{pmatrix}
V_1(x) u_1(x) \\ V_2(x) u_2(x) \\ \vdots \\ V_n(x) u_n(x)
\end{pmatrix}
, \quad u \in \HS.
\end{align*}

We denote the discrete Fourier transform by $\FS$;
\begin{align*}
\FS u(\xi) &=
\begin{pmatrix}
Fu_1(\xi) \\
Fu_2(\xi) \\
\vdots \\
Fu_n(\xi)
\end{pmatrix}, \quad \xi \in \T^d:=[-\pi,\pi)^d, \\
F u_j(\x) &= (2\pi)^{-\frac{d}{2}}\sum_{x\,\in\,\Z^d}e^{-ix\cdot\x}u_j(x),
\end{align*}
for $u\in \ell^1(\Z^d;\C^n)$. Then $\FS$ is extended to a unitary operator from $\HS$ onto $\hat{\HS}=L^2(\T^d; \C^n)$, and we denote its extension by the same symbol $\FS$. We easily see that $\FS \circ H_0 \circ \FS^*$ is the multiplication operator on $\T^d$ by the matrix-valued function
\begin{align*}
H_0(\xi) = 
\begin{pmatrix}
h_{11}(\xi) & h_{12}(\xi) & \cdots & h_{1n}(\xi) \\
h_{21}(\xi) & h_{22}(\xi) & \cdots & h_{2n}(\xi) \\
\vdots & \vdots & \ddots & \vdots \\
h_{n1}(\xi) & h_{n2}(\xi) & \cdots & h_{nn}(\xi)
\end{pmatrix},
\end{align*}
where
\begin{align*}
h_{jk}(\xi):=\sum_{x\,\in\,\Z^d}e^{-ix\cdot\xi}f_{jk}(x).
\end{align*}
Since $f_{jk}$'s are assumed to be rapidly decreasing, $h_{jk}$'s are smooth functions on $\T^d$. Note that $\sigma(H_0) = \{ \lambda \mid \det (H_0(\xi)-\lambda) =0 \ \text{for some}\ \xi \in \T^d \} $ and $H_0$ is a self-adjoint operator if and only if $H_0(\x)$ is a symmetric matrix for any $\x\in\T^d$, i.e., by the definition of $H_0(\x)$,
\begin{align}\label{selfadj}
\overline{f_{jk}(-x)}=f_{kj}(x), \quad x \in \Z^d, \ 1 \leq j,k \leq n.
\end{align}


In this paper, we assume the following assumption concerning the self-adjointness of $H_0$ and a long-range condition of $V$.

\begin{enonce}{Assumption}\label{ass}\leavevmode
\begin{enumerate}
\item\label{assum1.1.1} $f_{jk}$'s are rapidly decreasing functions satisfying~\eqref{selfadj}.

\item\label{assum1.1.2} $V={}^t(V_1, \cdots, V_n)$ has the following representation
\begin{align*}
V=V_L + V_S,
\end{align*}
where each entry of $V_L$ is the same, i.e., $V_L={}^t(V_\ell, \cdots, V_\ell)$ with some $V_\ell : \Z^d \to \R$. Furthermore, there exist $\rho>0$ and $C,C_\a>0$ such that
\begin{align}
\left|\tilde\partial_x^\a V_\ell(x)\right| &\leq C_\a \la x \ra^{-\rho-|\a|}, \label{long range}
\\
|V_S(x)| &\leq C \la x \ra^{-1-\rho} \label{short range}
\end{align}
for any $x\in\Z^d$ and $\a\in\Z_+^d$. Here $\tilde\partial_x^\a=\tilde\partial_{x_1}^{\a_1} \cdots \tilde\partial_{x_d}^{\a_d}$, $\tilde \partial_{x_j} V(x)=V(x)-V(x-e_j)$ is the difference operator with respect to the $j^{\rm th}$ variable.
\end{enumerate}
\end{enonce}


Assumption~\ref{ass} implies that $V$ is a compact operator on $\HS$ and hence
\begin{align*}
\sigma_{\operatorname{ess}}(H)=\sigma_{\operatorname{ess}}(H_0),
\end{align*}
where $\sigma_{\operatorname{ess}}(H)$ (resp.\ $\sigma_{\operatorname{ess}}(H_0)$) denotes the essential spectrum of $H$ (resp.\ $H_0$).


We denote the union of Fermi surfaces corresponding to the energies in $\Gamma\subset\R$ by
\begin{align*}
\Ferm(\Gamma):=&\,\left\{ p = (\x,\l)\in\T^d\times\Gamma \,\middle|\, \l \text{ is an eigenvalue of }H_0(\x) \right\} \\
=&\, \left\{ p = (\x,\l)\in\T^d\times\Gamma \,\middle|\, \det \left(H_0(\x)-\l\right)=0 \right\}. 
\end{align*}
Before describing the main theorem, we prepare the notation of non-threshold energies.

\begin{defi}\label{nonthr}
$\l_0\in\s(H_0)$ is said to be \emph{a non-threshold energy of $H_0$} if the following properties~\eqref{defi1.2.1} and~\eqref{defi1.2.2} hold:
\begin{enumerate}
\item\label{defi1.2.1} For any $\x_0\in\T^d$ such that $\det(H_0(\x_0)-\l_0)=0$, there exists an open neighborhood $G \subset \T^d \times \R$ of $p=(\x_0,\l_0)$ such that $\Ferm(\R) \cap G$ has a graph representation, i.e.
\begin{align}\label{lambda}
\Ferm(\R) \cap G = \{(\x,\l(\x)) \mid \x\in U \}
\end{align}
with some $U\ni \x_0$ and $\l\in C^\infty(U)$.

\item\label{defi1.2.2} Let $\x_0$ be arbitrarily fixed so that $\det(H_0(\x_0)-\l_0)=0$ holds, and let $\l(\x)$ be as in~\eqref{lambda}. Then $\nabla_\x \l(\x_0) \neq 0$ holds.
\end{enumerate}
\end{defi}

\begin{rema}
There is a sufficient condition of non-threshold energies:
\begin{align*}
\nabla_\x \det(H_0(\x)-\l_0)\neq0 \ \text{for any } \x\in\T^d \ \text{such that} \ \det(H_0(\x)-\l_0)=0,
\end{align*}
see Condition~(A-3) in~\cite[Sections~6~and~7]{An-I-M}. The principal difference is that Definition~\ref{nonthr} covers the case where $H_0(\x)$ has degenerate eigenvalues but no branching occurs.

On the other hand, the set of non-threshold energies in the present paper can be smaller than that of~\cite[Definition~5.5]{P-R}. Definition~\ref{nonthr} is assumed in order for each eigenvalue and its associated eigenprojection of the Fourier symbol $H_0(\x)$ of $H_0$ to be smooth with respect to $\x$ on the non-threshold energy level, which is needed to construct modifiers and to show some lemmas by the pseudodifferential calculus.
\end{rema}

Let $\Gamma(H_0)$ be the set of non-threshold energies of $H_0$. Then $\Gamma(H_0)$ is an open set of $\R$ and $\Gamma(H_0) \subset \s(H_0)$. Note that $H_0$ has purely absolutely continuous spectrum on $\Gamma(H_0)$, i.e., $\s_{pp}(H_0)\cap \Gamma(H_0)=\s_{sc}(H_0)\cap \Gamma(H_0)=\phi$ (see Remark~\ref{rem purely ac}).

The main theorem of this paper is the following.

\begin{theo}\label{main}
Suppose Assumption~\ref{ass} and $\Gamma \Subset \Gamma(H_0)$. Then there are bounded operators $J_\pm=J_{\pm,\Gamma}$ on $\HS$, called Isozaki--Kitada modifiers, such that the modified wave operators exist:
\begin{align}\label{modified WOs}
W_{\IK}^\pm(\Gamma)=\slim e^{itH}J_\pm e^{-itH_0}E_{H_0}(\Gamma),
\end{align}
where $E_{H_0}$ denotes the spectral measure of $H_0$, and that the following properties hold:
\begin{enumerate}\romanenumi
\item Intertwining property: $H W_{\IK}^\pm(\Gamma)=W_{\IK}^\pm(\Gamma) H_0$.
\item Partial isometries: $\|W_{\IK}^\pm(\Gamma) u\|=\|E_{H_0}(\Gamma)u\|$.
\item Completeness: $\Ran W_{\IK}^\pm(\Gamma) = E_{H}(\Gamma)\HS_{\ac}(H)$.
\end{enumerate}
\noindent Here $\HS_{\ac}(H)$ denotes the absolutely continuous subspace of $H$.
\end{theo}

Various examples of unperturbed operators $H_0$ are given by Ando, Isozaki and Morioka~\cite[Section~3]{An-I-M}. Note that, if the perturbation $V$ is short-range, i.e., $V_L=0$, we can set $J_\pm=\Idd_\HS$, thus there exist the wave operators in this case. See~\cite{P-R} for short-range scattering theory for discrete \Schr operators on various lattices. We also note that a long-range scattering theory in the case of $n=1$, e.g., discrete \Schr operators on square and triangular lattices, is considered by Nakamura~\cite{N} and the author~\cite{T1}. Moreover, Theorem~\ref{main} covers an arbitrary periodic lattice $\LS$ with each primitive unit cell $\LS/\Gamma$ containing finite elements, where $\Gamma \cong \Z^d$ denotes the transformation group associated to $\LS$. In particular, it includes the result by the author~\cite{T2}, where a long-range scattering theory for discrete \Schr operators on the hexagonal lattice is studied. See also~\cite{D-G,R-S3,Y} and references therein for scattering theory of \Schr operators on $\R^d$.



The key idea of the present paper is to locally observe the eigenvalues $\lambda_{k}(\x)$ of $H_0(\x)$. This enables us to construct modifiers via the Hamilton flow on $T^*\T^d=\R^d \times \T^d$ associated with $\lambda_k(\x)+\tilde{V}_{\ell}(x)$, where $\tilde{V}_{\ell}$ is a smooth extension of $V_\ell$ onto $\R^d$ satisfying $|\partial_x^\a \tilde V_\ell(x)|\leq C_\a' \la x \ra^{-\rho-|\a|}$, as well as to have the limiting absorption principle and the radiation estimate. For local observation of eigenvalues, we need to use the eigenprojections of $H_0(\x)$ depending smoothly on $\x$. We note that the concrete diagonalization of $H_0(\x)$ is employed in~\cite{T2}, where the representation of diagonalization is smoothly defined globally away from the Dirac points.

The organization of this paper is as follows. We first prepare notations and properties of pseudodifference operators in Section~\ref{notations sec}. In Section~\ref{Mourre section}, the limiting absorption principle and the propagation estimate for $H$ are studied. We use the Mourre theory and a standard argument of the propagation of wave packets as in Yafaev~\cite[Chapter~10]{Y}. The construction of conjugate operators is essentially due to Parra and Richard~\cite{P-R}. Section~\ref{classical section} is devoted to constructing phase functions which are given as local solutions to eikonal equations corresponding to each fiber of eigenvalues of $H_0(\x)$. The construction of phase functions is due to~\cite{N3}. In Section~\ref{IK section}, using the phase functions in the previous section, we construct Isozaki--Kitada modifiers. Finally in Section~\ref{proof section}, we use lemmas in the previous section to prove Theorem~\ref{main}. The proof is based on Kato's smooth perturbation theory, and is an analogue of that in long-range scattering theory for \Schr operators on $\R^d$ (see~\cite{Y}).



\section{Preliminaries}\label{notations sec}



\subsection{Representations of fibers}

Let $\Gamma$ be as in Theorem~\ref{main}, and let $I \Subset \Gamma(H_0)$ be fixed so that $\Gamma \Subset I \Subset \Gamma(H_0)$.

For each $p=(\x_0,\l_0) \in \Ferm(\Gamma(H_0))$, let $G=G_p$ be as in Definition~\ref{nonthr}. Then $\{ G_p \}_{p\,\in\,\Ferm(\Gamma(H_0))}$ is an open covering of $\Ferm(\Gamma(H_0))$. Since $\Ferm(\overline{I})$ is compact, we can take a finite family $\{ G_{j} \}_{j=1}^J = \{ G_{p_j} \}_{j=1}^J$ of open sets which covers $\Ferm(\overline{I})$.

Note that $\{ G_j \cap \Ferm(\R) \}_{j=1}^J$ is also a covering family of $\Ferm(\overline{I})$. Let $G'_k$, $k=1,\dots,K$, be the connected components of $\cup_{j=1}^J G_j \cap \Ferm(\R)$. We see that each $G'_k$ remains to have a graph representation
\begin{align}\label{G_k representation}
G'_k=\left\{ (\x,\l_k(\x)) \,\middle|\, \x \in \US_k \right\}
\end{align}
with some open set $\US_k \subset \T^d$ and $\l_k\in C^\infty(\US_k)$. We denote by $P_k(\x)$ the projection matrix onto $\kerr (H_0(\x)-\l_k(\x))$ for $\x\in \US_k$. Then we have for $\psi\in C_c^\infty(I)$
\begin{align}\label{Rep of E_H0}
\ps(H_0(\x)) = \sum_{k=1}^K \ps(\l_k(\x)) P_k(\x) \chi_{\US_k}(\x).
\end{align}


\subsection{Pseudodifference calculus}

For $a:\Z^d \times \T^d \to M_n(\C)\cong \C^{n\times n}$,
\begin{align*}
a(x,D_x)u(x) :=(2\pi)^{-\frac{d}{2}}\int_{\T^d} e^{ix\cdot\x} a(x,\x) \FS u(\x)d\x, \quad u\in \HS,
\end{align*}
denotes the pseudodifference operator on $\Z^d$ with symbol $a(x,\x)$. If $a$ depends only on $\x$, we denote by $a(D_x)=\FS^*\circ a(\cdot) \circ \FS$ the Fourier multiplier associated with $a(\x)$ in short.

We cite a lemma concerning the pseudodifference calculus on $\HS$ (see~\cite[Theorem~4.2.10]{R-T} and the proof of~\cite[Lemma~2.2]{T1}).

\begin{lemm}\label{PDO lemma}
Let $a:\Z^d\times\Z^d\times\T^d \to M_n(\C)$ be a smooth function with respect to $\T^d$, and let
\begin{align*}
Au(x) = (2\pi)^{-d} \int_{\T^d} \sum_{y\,\in\,\Z^d} e^{i(x-y)\cdot\x} a(x,y,\x) u(y)d\x.
\end{align*}
Suppose that for any $\a\in\Z_+^d$
\begin{align}\label{S^0}
\sup_{(x,y,\x)\,\in\,\Z^d\times\Z^d\times\T^d}\left|\partial_{\x}^\a a(x,y,\x)\right| < \infty.
\end{align}
Then $A$ is a bounded operator on $\ell^2(\HS)$.
\end{lemm}

Let $S^m$ be the symbol class of order $m\in\R$, i.e.,
\begin{align*}
S^m=\left\{ a: \Z^d\times\T^d \to M_n(\C) \;\middle|\;
\begin{aligned}
& a(x,\cdot)\in C^\infty\left(\T^d;M_n(\C)\right),\ \forall x\in\Z^d,\\ 
&\sup_{(x,\x)\,\in\,\Z^d\times\T^d} \jap{x}^{-m+|\a|} 
\left|\tilde\partial_{x}^\a\partial_{\x}^\b a(x,\x)\right|\\
&< \infty,\ \forall \a,\b\in\Z_+^d
\end{aligned}
\right\},
\end{align*}
where $\tilde\partial_{x}^\a$ denotes the difference operator as in~\eqref{long range}.

The following two assertions are analogous to the composition formula for pseudodifferential operators. See~\cite[Theorems~4.7.3 and 4.7.10]{R-T} for the proofs.

\begin{lemm}\label{composition lemma}
Let $a\in S^m$ and $b\in S^{\ell}$. Then $a(x,D_x)b(x,D_x)=c(x,D_x)$ with some $c\in S^{m+\ell}$ satisfying the asymptotic expansion
\begin{align*}
&c(x,\x) - \sum_{|\a|\,\leq\,M} \partial_{\x}^\a a(x,\x) \tilde\partial_{x}^\a b(x,\x) \in S^{m+\ell-M-1}
\end{align*}
for any $M\in\Z_+$.
\end{lemm}

\begin{lemm}\label{adjoint lemma}
Let $a\in S^m$. Then there exists $b\in S^m$ such that $a(x,D_x)^*=b(x,D_x)$ and $b(x,\x) - a(x,\x)^* \in S^{m-1}$.
\end{lemm}


\subsection{Kato's smooth perturbation theory} For a self-adjoint operator $H$ and an $H$-bounded operator $G$, we say that $G$ is $H$-smooth if
\begin{align}\label{Kato smooth}
\frac{1}{2\pi}\sup_{\|u\|_{\HS}=1,u\,\in\,D(H)} \int_{-\infty}^\infty \left\| Ge^{-itH}u \right\|^2 dt < \infty.
\end{align}
For a Borel set $I\subset\R$, we say that $G$ is $H$-smooth on $I$ if $GE_H(I)$ is $H$-smooth, and we also say that $G$ is locally $H$-smooth on $I$ if $G$ is $H$-smooth on $I'$ for any $I' \Subset I$.

There are several conditions equivalent to~\eqref{Kato smooth} (see e.g.~\cite{Yg}), and the one we need in the following is:
\begin{align}\label{Kato smooth 2}
\sup_{\lambda\,\in\,\R,\,\e\,>\,0} \left\|G\d_\e(\lambda,H)G^*\right\| <\infty,
\end{align}
where $\d_\e(\lambda,H)=\frac{1}{2\pi i}\{(H-\lambda-i\e)^{-1}-(H-\lambda+i\e)^{-1}\}$.


\section{Limiting absorption principle and radiation estimates}\label{Mourre section}

In this section, we consider the limiting absorption principle and radiation estimates for the proof of Theorem~\ref{main}.


\subsection{Limiting absorption principle}


For a self-adjoint operator $A$ and $m\in\bbN$, let
\[
C^m(A)=\left\{ S\in\BS(\HS) \,\middle|\, \R\to\BS(\HS),t\mapsto e^{-itA}Se^{itA} \ \text{is strongly of class} \ C^m \right\},
\]
and $C^\infty(A)=\cap_{m\,\in\,\bbN}C^m(A)$. We denote by $\CS^{1,1}(A)$ the set of the operators $S$ satisfying
\[
\int_0^1 \left\| e^{-itA} S e^{itA}+e^{itA} S e^{-itA}-2S \right\| \frac{dt}{t^2} < \infty.
\]


We set the Besov space
\begin{align*}
B:=(\DS(\jap{x}),\HS)_{\frac{1}{2},1},
\end{align*}
where we have used the notation of real interpolation $(\cdot,\cdot)_{\theta,p}$ between Banach spaces (see~\cite[Section~2.1]{A-BdM-G}).

The following proposition is called the limiting absorption principle. The proof is given by the Mourre theory, and the construction of conjugate operators is essentially due to~\cite[Lemma~6.2]{P-R}.

\begin{prop}\label{LAP prop}
Suppose Assumption~\ref{ass}. Then:
\begin{enumerate}
\item\label{prop3.1.1} The set of eigenvalues of $H$ is locally finite in $\Gamma(H_0)$ with counting multiplicities.

\item\label{prop3.1.2} For any $\lambda \in \Gamma(H_0)\backslash \s_{pp}(H)$, there exist the weak-* limits in $\BS(B,B^*)$
\begin{align*}
\operatorname{w^*-}\lim_{\e\to+0} (H-\l\mp i\e)^{-1}.
\end{align*}
Moreover, each convergence is locally uniform in $\l\in\Gamma(H_0)\backslash \s_{pp}(H)$. In particular, for any $\Gamma\Subset\Gamma(H_0)\backslash \s_{pp}(H)$,
\begin{align}\label{LAP}
\sup_{\l\,\in\,\Gamma,\,\e\,>\,0} \left\|(H-\l\mp i\e)^{-1}\right\|_{\BS(B,B^*)} < \infty.
\end{align}
\end{enumerate}
\end{prop}

\begin{proof}
Let $\Gamma\Subset\Gamma(H_0)$ be arbitrarily fixed, and recall the representation~\eqref{Rep of E_H0}. We set $\chi_k \in C_c^\infty(\US_k)$ so that $\chi_k=1$ on $\lambda_k^{-1}(\Gamma)$. We also set the conjugate operator $A$ by
\begin{align*}
A &= \sum_{k=1}^K P_k(D_x) \chi_k(D_x) i\left[\l_k(D_x),|x|^2\right] P_k(D_x) \chi_k(D_x) \\
&= \sum_{k=1}^K P_k(D_x) \chi_k(D_x) M_k P_k(D_x) \chi_k(D_x), 
\end{align*}
where
\begin{align*}
M_k = x \cdot \nabla_\x\l_k(D_x) + \nabla_\x\l_k(D_x) \cdot x.
\end{align*}
Now we employ the Mourre theory (\cite[Proposition~7.1.3, Corollary~7.2.11, Theorem~7.3.1]{A-BdM-G}, see also~\cite[Theorem~A.1]{T2}). Then, since $A$ is $\jap{x}$-bounded, it suffices to show that $H\in \CS^{1,1}(A)$ and that, for any $\psi \in C_c^\infty(\Gamma)$, there exist $c>0$ and a compact operator $K$ such that the Mourre inequality holds:
\begin{align}
\psi(H) i[H,A] \psi(H) \geq c \psi(H)^2 + K. \label{Mourre}
\end{align}

For the first assertion, we easily see $H_0\in C^\infty(A)$, and $V \in \CS^{1,1}(A)$ is proved by~\eqref{long range}, \eqref{short range} and Lemma~\ref{composition lemma} (see~\cite{P-R,T2} for details of the proof).

For the proof of~\eqref{Mourre}, we learn by Definition~\ref{nonthr}$\MK$\eqref{defi1.2.2} that
\begin{multline}\label{free Mourre}
\psi(H_0) i[H_0,A] \psi(H_0) \\
\begin{aligned}
&= 2 \sum_{k=1}^K P_k(D_x) \ps\bigl(\l_k(D_x)\bigr) \chi_k(D_x) |\nabla_\x \lambda_k(D_x)|^2 P_k(D_x) \ps\bigl(\l_k(D_x)\bigr) \chi_k(D_x) \\
&\geq c \sum_{k=1}^K P_k(D_x) \ps\bigl(\l_k(D_x)\bigr)^2 \chi_k(D_x)^2 \geq c \ps(H_0)^2. 
\end{aligned}
\end{multline}
It follows from~\eqref{long range} and~\eqref{short range} that $i[V,A]$ and $\psi(H)-\psi(H_0)$ are compact, and hence we have~\eqref{Mourre}.
\end{proof}

\begin{rema}\label{rem purely ac}
If we adopt the Mourre theory to $H=H_0$, \eqref{free Mourre} implies that $H_0$ has purely absolutely continuous spectrum on $\Gamma(H_0)$.
\end{rema}

Since $B \supset \jap{x}^s\HS$ and $B^* \subset \jap{x}^{-s}\HS$ hold for any $s>\frac{1}{2}$ (see, e.g., \cite[Theorem~3.4.1]{A-BdM-G}), \eqref{LAP} and the equivalence between~\eqref{Kato smooth} and~\eqref{Kato smooth 2} imply the following corollary.

\begin{cor}\label{cor_lap}
For any $s>\frac{1}{2}$, $\jap{x}^{-s}$ is locally $H$-smooth on $\Gamma(H_0) \backslash \s_\pp(H)$.
\end{cor}



\subsection{Radiation estimates} In order to prove the existence and completeness of modified wave operators, we use, in addition to the limiting absorption principle, other propagation estimates called radiation estimates (see~\cite[Theorem~10.1.7]{Y}).

\begin{prop}\label{prop_rad_est}
Let $\Gamma\Subset\Gamma(H_0)$ be fixed, and let $\l_k(\x)$, $k=1,\dots,K$, be as in~\eqref{G_k representation}. We set for $k=1,\dots,K$ and $j=1,\dots,d$,
\begin{align*}
\nabla_{k,j}^\perp := \left\{\left(\partial_{\x_j}\l_k\right)(D_x) - \chi_{\{x\neq0\}} |x|^{-2} x_j \la x, (\nabla_\x \l_k)(D_x)\ra \right\} P_k(D_x)\chi_k(D_x),
\end{align*}
where $\chi_k \in C_c^\infty(\US_k)$ is fixed arbitrarily so that $\chi_k=1$ on $\lambda_k^{-1}(\Gamma)$. Then
\begin{align}\label{radiation op}
\chi_{\{x\neq0\}} |x|^{-\frac{1}{2}} \nabla_{k,j}^\perp
\end{align}
is locally $H$-smooth on $\Gamma(H_0)\backslash \s_{pp}(H)$.
\end{prop}

\begin{proof}
Fix $k=1,\dots,K$. For simplicity of notation, we write $\l$, $P$, $\chi$ and $\nabla_j^\perp$ instead of $\l_k$, $P_k$, $\chi_k$ and $\nabla_{k,j}^\perp$, respectively.

Let $a\in C^\infty(\R^d)$ be fixed so that $a(x)=|x|$ for $|x|\geq1$, and let
\begin{align*}
a_j:=\partial_{x_j}a, \quad v_j:=\partial_{\x_j}\l.
\end{align*}
We set
\begin{align*}
\bbA := (P\chi)(D_x) \sum_{j=1}^d \left\{a_j(x)v_j(D_x)+v_j(D_x)a_j(x) \right\} (P\chi)(D_x).
\end{align*}
Then $\bbA$ is a bounded symmetric operator. It follows from~\cite[Proposition~0.5.11]{Y} and Corollary~\ref{cor_lap} that we only have to show
\begin{align}\label{radiation ineq}
\left(u, i[H,\bbA]u\right)
\geq 2\sum_{j=1}^d \left\| \chi_{\{x\neq0\}}|x|^{-1/2} \nabla_{j}^\perp u \right\|^2 -C\|\jap{x}^{-\mu}u\|^2
\end{align}
with some $C>0$, where $\mu=\min(\frac{\rho+1}{2},1)>\frac12$.

We show~\eqref{radiation ineq} in the rest of the proof. The representation~\eqref{Rep of E_H0} implies
\begin{align*}
i[H_0,\bbA] = (P\chi)(D_x) \cdot M \cdot (P\chi)(D_x),
\end{align*}
where
\begin{align*}
M = \sum_{j=1}^d \left\{i\left[\l(D_x),a_j(x)\right] \cdot v_j(D_x) + v_j(D_x) \cdot i\left[\l(D_x),a_j(x)\right] \right\}.
\end{align*}
It follows from Lemma~\ref{composition lemma} that, formally,
\begin{align*}
M &= 2 \sum_{j=1}^d \sum_{\ell=1}^d v_\ell(D_x)a_{j\ell}(x) v_j(D_x) + R_1,
\end{align*}
where $a_{j\ell}:=\partial_{x_\ell}\partial_{x_j}a$, and $R_1$ satisfies $\jap{x}^2 (P\chi)(D_x) R_1 (P\chi)(D_x) \in \BS(\HS)$. Since for $|x|\geq1$
\begin{align*}
a_{j\ell}(x) = \partial_{x_\ell}\partial_{x_j}\left(|x|\right) = -\frac{x_j x_\ell}{|x|^3} + \d_{j\ell} |x|^{-1},
\end{align*}
we learn
\begin{equation}\label{radiation pf commutator}
\begin{split}
(u,i[H_0,\bbA]u) 
&= -2 \sum_{j=1}^d \sum_{\ell=1}^d \left(u^\ell, \frac{x_j x_\ell}{|x|^3} \chi_{\{x\neq0\}} u^j\right) \\
&\quad+2 \sum_{j=1}^d \left(u^j, |x|^{-1}\chi_{\{x\neq0\}}u^j\right)
 + \bigl((P\chi)(D_x)u, R_2 (P\chi)(D_x)u\bigr), 
\end{split}
\end{equation}
where
\begin{align*}
u^j:= (v_j P\chi)(D_x) u,
\end{align*}
and
\begin{align*}
R_2= R_1 + 2 \sum_{j=1}^d \sum_{\ell=1}^d a_{j\ell}(0) v_\ell(D_x) \chi_{x=0}(x) v_j(D_x)
\end{align*}
also satisfies $\jap{x}^2 (P\chi)(D_x) R_2 (P\chi)(D_x) \in \BS(\HS)$.

On the other hand, a direct computation implies for $x\neq0$
\begin{multline*}
\left| \nabla_{j}^\perp u (x) \right|^2 
= \left|u^j(x)\right|^2 - |x|^{-2} x_j \sum_{\ell=1}^d x_\ell \left(u^\ell(x)\overline{u^j(x)}+\overline{u^\ell(x)}u^j(x) \right) \\
+ |x|^{-4} {x_j}^2 \sum_{\ell=1}^d \sum_{m=1}^d x_\ell x_m \overline{u^\ell(x)}u^m(x).
\end{multline*}
Summing up over $j=1,\dots,d$, we learn
\begin{equation}\label{differentiation across trajectory}
\begin{split}
\sum_{j=1}^d \left| \nabla_{j}^\perp u (x) \right|^2 
&= \sum_{j=1}^d \left|u^j(x)\right|^2 - |x|^{-2} \sum_{j=1}^d \sum_{\ell=1}^d x_j x_\ell \left(u^\ell(x)\overline{u^j(x)}+\overline{u^\ell(x)}u^j(x) \right) \\
& \qquad + |x|^{-2} \sum_{\ell=1}^d \sum_{m=1}^d x_\ell x_m \overline{u^\ell(x)}u^m(x) \\
&= \sum_{j=1}^d \left|u^j(x)\right|^2 - |x|^{-2} \sum_{\ell=1}^d \sum_{m=1}^d x_\ell x_m \overline{u^\ell(x)}u^m(x), \quad x\neq0.
\end{split}
\end{equation}

Combining~\eqref{differentiation across trajectory} with~\eqref{radiation pf commutator}, we obtain
\begin{multline*}
\left(u, i[H,\bbA]u\right)\\
= 2\sum_{j=1}^d \left\| \chi_{\{x\neq0\}}|x|^{-1/2} \nabla_{j}^\perp u \right\|^2 
+ \bigl((P\chi)(D_x)u, R_2 (P\chi)(D_x)u\bigr) + (u, i[V,\bbA] u).
\end{multline*}
We see that $\jap{x}^{1+\rho} [V,\bbA] \in \BS(\HS)$ by~\eqref{long range}, \eqref{short range} and Lemma~\ref{composition lemma}. Finally we obtain~\eqref{radiation ineq}.
\end{proof}


\section{Classical mechanics}\label{classical section}

In this section, we construct phase functions used for the definition of time-independent modifiers $J_\pm$ in~\eqref{modified WOs}. For the precise definition of $J_\pm$, see~\eqref{time-independent modifiers}.

Let $\l_k(\x):\US_k\to\R$, $k=1,\dots,K$, be the functions in~\eqref{G_k representation}. The next proposition concerns the classical scattering problem with respect to the Hamiltonian $\l_k(\x)+\tilde V_\ell(x)$ on $T^*\US_k=\R_x^d\times\US_k$, where $\tilde V_\ell$ is a smooth extension of $V_\ell$ onto $\R^d$ such that $|\partial_x^\a \tilde V_\ell(x)|\leq C_\a' \la x \ra^{-\rho-|\a|}$ holds. See~\cite[Lemma~2.1]{N} for a concrete construction of $\tilde V_\ell$.

The proof of the following proposition is given by~\cite[Section~2]{N3} (see also~\cite{IsoKita,T1}).

\begin{prop}\label{phase functions}
Let $\l_k(\x):\US_k\to\R$, $k=1,\dots,K$, be fixed. Then for any open set $U \Subset \US_k$ and $\e\in(0,2)$, there exist $R>0$ and smooth functions $\f_\pm^k(x,\x)$ defined on a neighborhood of
\begin{align*}
D_{k,\pm} = \left\{(x,\x)\in \R^d \times U \,\middle|\, |x| \geq R, \ \pm \cos(x,\nabla\lambda_k(\x)) \geq-1+\e \right\},
\end{align*}
where
\begin{align*}
\cos(x,\nabla\l_k(\x)) := \frac{x\cdot\nabla\l_k(\x)}{|x| |\nabla\l_k(\x)|},
\end{align*}
such that
\begin{align}\label{eikonal}
\l_k\left(\nabla_x\f_\pm^k(x,\x)\right) + \tilde V_\ell(x) = \l_k(\x), \quad (x,\x) \in D_{k,\pm}.
\end{align}
Furthermore, $\f_\pm^k$ satisfy for $(x,\xi)\in D_{k,\pm}$
\begin{align}
\left|\partial_x^\alpha\partial_\xi^\beta\left[\varphi_\pm^k(x,\xi)-x\cdot\xi\right]\right| &\leq C_{\alpha\beta} \langle x\rangle^{1-\rho-|\alpha|}, \label{phase estimate}
\\
\left|{}^t\nabla_x\nabla_\xi\varphi_\pm^k(x,\xi)-I \right|&<\frac{1}{2}. \label{phase estimate2}
\end{align}
\end{prop}


\section{Construction of Isozaki--Kitada modifiers}\label{IK section}

Let $\Gamma\Subset\Gamma(H_0)$ be fixed. Let $\l_k \in C^\infty(\US_k)$, $k=1,\dots,K$, be as in~\eqref{G_k representation}, and let $\f_{\pm}^{k}$ be the phase functions constructed in Proposition~\ref{phase functions} with setting $\e=\frac14$ and $U$ so that $\l_k^{-1}(\Gamma)\Subset U \Subset \US_k$.

We take functions $\chi_k\in C_c^\infty(U;[0,1])$, $\eta\in C^\infty(\R^d)$ and $\s_\pm\in C^\infty(\R;[0,1])$ such that
\begin{align}
&\chi_k(\x)=1, \quad \x \in \lambda_k^{-1}(\Gamma), \label{chi 1}
\\
&\eta(x)=
\begin{alcases}
1 & \quad \text{if } |x|\geq 2R, \\
0 & \quad \text{if } |x|\leq R,
\end{alcases}\label{eta 01}
\\
&\s_\pm(\theta)=
\begin{alcases}
1 & \quad \text{if } \pm\theta \geq \frac{1}{2}, \\ 0 & \quad \text{if } \pm\theta \leq -\frac{1}{2},
\end{alcases}\label{sigma 01}
\\
&\s_+(\theta)^2+\s_-(\theta)^2=1, \quad \theta\in\R, \label{sigma sum1}
\end{align}
where $R>0$ is the constant in Proposition~\ref{phase functions}. Then we define the Isozaki--Kitada modifiers $J_\pm^k$ associated with the pair $(P_k,\l_k,\US_k)$ by
\begin{align}\label{IK1}
J_\pm^k u(x) := (2\pi)^{-\frac{d}{2}} \int_{\T^d} e^{i\f_{\pm}^k(x,\x)} s_\pm^k(x,\x) \FS u(\x) d\x,
\end{align}
where
\begin{align*}
s_\pm^k(x,\x):=\eta(x)\s_\pm\bigl(\cos(x,\nabla\l_k(\x)) \bigr) P_k(\x) \chi_k(\x).
\end{align*}
We recall that $P_k(\x)$ is the projection matrix onto $\kerr(H_0(\x)-\l_k(\x))$, and note that $\supp s_\pm^k \subset D_{k,\pm}$ holds. Their formal adjoints are given by
\begin{align*}
\left(J_\pm^k\right)^* u(x) = \FS^* \left((2\pi)^{-\frac{d}{2}} \sum_{y\,\in\,\Z^d} e^{-i\f_{\pm}^k(y,\cdot)} s_\pm^k(y,\cdot) u(y) \right).
\end{align*}
Direct computations imply
\begin{align}\label{s is in S^0}
\sup_{(x,\x)\,\in\,\R^d\times\T^d}\jap{x}^{|\a|}\left|\partial_x^\a\partial_{\x}^\b s_\pm^k(x,\x)\right| < \infty,
\end{align}
in particular~\eqref{S^0} holds.

The next lemma follows from an analogue of the argument of calculus of Fourier integral operators (see~\cite{IsoKita,T1}).

\begin{lemm}\label{J lem}
Let $k=1,\dots,K$ be fixed, and let $\rho>0$ be the constant in Assumption~\ref{ass}$\MK$\eqref{assum1.1.2}. Then:
\begin{enumerate}
\item\label{lemm5.1.1} $J_\pm^k$ are bounded operators on $\HS$.

\item\label{lemm5.1.2} The operators
\begin{align}
&\la x \ra^\rho \left(J_\pm^k \left(J_\pm^k\right)^* - s_\pm^k(x,D_x) s_\pm^k(x,D_x)^* \right), \label{J2}
\\
&\la x \ra^\rho \left(\left(J_\pm^k\right)^* J_\pm^k - s_\pm^k(x,D_x)^* s_\pm^k(x,D_x) \right) \label{J3}
\end{align}
are bounded on $\HS$.
\item\label{lemm5.1.3} For any $q\geq0$,
\begin{align}
\la x \ra^{-q} J_\pm^k \la x \ra^q, \label{JX}
\end{align}
is bounded on $\HS$.
\item\label{lemm5.1.4} Suppose that $\psi = \psi(\x) \in C^\infty(\T^d;M_n(\C))$ commutes with $s_\pm^k(x,\x)$ for any $(x,\x)\in\Z^d\times\T^d$. Then
\begin{align}
\la x \ra^\rho \left[J_\pm^k, \psi(D_x)\right] \label{JXI}
\end{align}
is bounded on $\HS$. In particular, $[J_\pm^k, \psi(D_x)]$ are compact.
\item\label{lemm5.1.5} If $k \neq \ell$, then $J_\pm^k (J_\pm^\ell)^*=0$, and $(J_\pm^k)^* J_\pm^\ell$ are compact on $\HS$.
\end{enumerate}
\end{lemm}


\begin{proof}
\eqref{lemm5.1.1} We compute
\begin{align*}
J_\pm^k \left(J_\pm^k\right)^* u(x) = (2\pi)^{-d} \int_{\T^d} \sum_{y\,\in\,\Z^d} e^{i\left(\f_{\pm}^k(x,\x)-\f_{\pm}^k(y,\x)\right)} s_\pm^k(x,\x) s_\pm^k(y,\x) u(y) d\x.
\end{align*}
We set $\f_{\pm}^k(x,\x)-\f_{\pm}^k(y,\x) = (x-y) \cdot \z(\x;x,y)$, where
\begin{align*}
\z(\x;x,y) := \int_0^1 \nabla_x \f_{\pm}^k(y+\theta(x-y),\x) d\theta.
\end{align*}
Then Proposition~\ref{phase functions} implies that the mapping $\x \mapsto \z(\x;x,y)$ is a diffeomorphism from $U$ into $\z(U)$ for any $x,y \in \Z^d$. Thus we have
\begin{align*}
J_\pm^k \left(J_\pm^k\right)^* u(x) = (2\pi)^{-d} \int_{\T^d} \sum_{y\,\in\,\Z^d} e^{i(x-y)\cdot\z} t_\pm^k(x,y,\z) u(y) d\z,
\end{align*}
where
\begin{align*}
t_\pm^k(x,y,\z):=s_\pm^k(x,\x(\z;x,y)) s_\pm^k(y,\x(\z;x,y)) \left|\det\left(\frac{d\x}{d\z} \right)\right|.
\end{align*}
Since $|\frac{d\z}{d\x}(\x) - I| < \frac{1}{2}$ by Proposition~\ref{phase functions}, \eqref{s is in S^0} implies $|\partial_{\z}^\a t_\pm^k(x,y,\z)|\leq C_\a$ for any $\a$. Therefore $J_\pm^k$ are bounded by Lemma~\ref{PDO lemma}.


\eqref{lemm5.1.2} The same argument as in~\eqref{lemm5.1.1} implies
\begin{align*}
\left(J_\pm^k \left(J_\pm^k\right)^* - s_\pm^k(x,D_x) s_\pm^k(x,D_x)^* \right) u(x) 
= (2\pi)^{-d} \int_{\T^d} \sum_{y\,\in\,\Z^d} e^{i(x-y)\cdot\z} r(x,y,\z) u(y) d\z,
\end{align*}
where
\begin{align*}
r(x,y,\z)=t_\pm^k(x,y,\z) - s_\pm^k(x,\z) s_\pm^k(y,\z).
\end{align*}
Since $|\partial_\x^\a r(x,\z,y)| \leq C_\a' \la x \ra^{-\rho}$, Lemma~\ref{PDO lemma} implies the boundedness of~\eqref{J2}.

The other case~\eqref{J3} can be treated similarly if we consider the justification of PDO calculus; the argument using Poisson's summation formula as in~\cite[Lemma~7.1]{N3} (see also~\cite[Lemma~2.3]{T1}) implies
\begin{multline*}
\FS \left(J_\pm^k\right)^* J_\pm^k \FS^* f(\x) \\
\begin{aligned}
=&\,(2\pi)^{-d} \int_{\R^d} \int_{\T^d} e^{i\left(-\f_{\pm}^k(x,\x)+\f_{\pm}^k(x,\eta)\right)} s_\pm^k(x,\x) s_\pm^k(x,\eta) f(\eta) d\eta dx + K_1 f(\x), \\
\FS s_\pm^k(x,D_x)^* & s_\pm^k(x,D_x) \FS^* f(\x) \\
=&\,(2\pi)^{-d} \int_{\R^d} \int_{\T^d} e^{i x\cdot(-\x+\eta)} s_\pm^k(x,\x) s_\pm^k(x,\eta) f(\eta) d\eta dx + K_2 f(\x),
\end{aligned}
\end{multline*}
where $K_j$, $j=1,2$, is a smoothing operator in the sense that $\jap{D_x}^N K_j \in \BS(\HS)$ for any $N>0$. Then by changing variables $x\mapsto \int_0^1 \nabla_\x\f_\pm^k(x,\x+\theta(\eta-\x))d\theta$, PDO calculus on $\T^d$ implies the boundedness of~\eqref{J3}.

\eqref{lemm5.1.3} By a complex interpolation argument, it suffices to show~\eqref{JX} for $q\in 2\Z_+$. Note that for $\a\in\Z_+^d$
\begin{align*}
 J_\pm^k x^\a u(x) 
&= (2\pi)^{-\frac{d}{2}} \int_{\T^d} e^{i\f_{\pm}^k(x,\x)}s_\pm^k(x,\x) i^{|\a|}\partial_{\x}^\a \FS u(\x) d\x \\
&= (-i)^{|\a|} (2\pi)^{-\frac{d}{2}} \int_{\T^d} \partial_{\x}^\a \Bigl(e^{i\f_{\pm}^k(x,\x)} s_\pm^k(x,\x) \Bigr) \FS u(\x) d\x. 
\end{align*}
Then we learn for any $N\in\Z_+$,
\begin{align*}
J_\pm^k \jap{x}^{2N} u(x) = (2\pi)^{-\frac{d}{2}} \int_{\T^d} e^{i\f_{\pm}^k(x,\x)} \left(L^{N}s_\pm^k\right)(x,\x) \FS u(\x) d\x,
\end{align*}
where $L := \jap{\nabla_\x\f_\pm^k}^2 - i\Delta_\x\f_\pm^k - 2i\jap{\nabla_\x\f_\pm^k,\nabla_\x} - \Delta_\x$. Since
\begin{align*}
\left|\partial_\x^\b\left(L^{N}s_\pm^k\right)(x,\x)\right| \leq C_{p,\b,N} \jap{x}^{2N}
\end{align*}
for any $\b\in\Z_+^d$, we have the boundedness of~\eqref{JX}.



\eqref{lemm5.1.4} It suffices to show the boundedness of $\jap{D_\x}^\rho [\hat{J}_\pm^k, \psi(\x)]$ as an operator on $L^2(\T^d;\C^n)$, where $\hat{J}_\pm^k:=\FS J_\pm^k\FS^*$. A direct computation implies
\begin{multline*}
\la D_\x \ra^\rho \left[\hat{J}_\pm^k, \psi(\x)\right]f(\x) \\
\begin{aligned}
&= (2\pi)^{-d} \sum_{x\,\in\,\Z^d}\int_{\T^d} e^{i\left(-x\cdot\x+\f_\pm^k(x,\y)\right)} \jap{x}^\rho \bigl(\psi(\y)-\psi(\x)\bigr) s_\pm^k(x,\y) f(\y) d\y \\
&= (2\pi)^{-d} \sum_{x\,\in\,\Z^d}\int_{\T^d} e^{i\left(-x\cdot\x+\f_\pm^k(x,\y)\right)} \jap{x}^\rho \Psi_1(x,\y) s_\pm^k(x,\y) f(\y) d\y  \\
&\quad + (2\pi)^{-d} \sum_{x\,\in\,\Z^d}\int_{\T^d} e^{i\left(-x\cdot\x+\f_\pm^k(x,\y)\right)} \jap{x}^\rho \Psi_2(x,\x,\y) s_\pm^k(x,\y) f(\y) d\y, 
\end{aligned}
\end{multline*}
where
\begin{align*}
\Psi_1(x,\y)\, &{:=}\,  \psi(\y)-\psi\left(\nabla_x\f_\pm^k(x,\y)\right), \\
\Psi_2(x,\x,\y)\, &{:=}\,  \psi\left(\nabla_x\f_\pm^k(x,\y)\right)-\psi(\x).
\end{align*}
The first term is treated similarly to~\eqref{lemm5.1.2}, since $|\partial_\y^\a \Psi_1(x,\y)| \leq C_\a \jap{x}^{-\rho}$ by~\eqref{phase estimate}. For the second term, we first employ the argument in the proof of boundedness of~\eqref{J3} to replace the summation over $\Z^d$ by the integral on $\R^d$ modulo smoothing operators. Then, since
\begin{align*}
\Psi_2(x,\x,\y) = \left(\nabla_x\f_\pm^k(x,\y)-\x\right)\cdot\int_0^1 \nabla_\x \psi\left(\x+\theta\left(\nabla_x\f_\pm^k(x,\y)-\x\right)\right) d\theta,
\end{align*}
we have
\begin{multline*}
(2\pi)^{-d} \int_{\R^d}\int_{\T^d} e^{i\left(-x\cdot\x+\f_\pm^k(x,\y)\right)} \jap{x}^\rho \Psi_2(x,\x,\y) s_\pm^k(x,\y) f(\eta) d\eta dx \\
= i(2\pi)^{-d} \int_{\R^d}\int_{\T^d} e^{i\left(-x\cdot\x+\f_\pm^k(x,\y)\right)} a(\x,\y,x) f(\y) d\y dx,
\end{multline*}
where
\begin{align*}
a(\x,\y,x)=\nabla_x \cdot \left(\jap{x}^\rho s_\pm^k(x,\y) \int_0^1 \nabla_\x\psi\left(\x+\theta\left(\nabla_x\f_\pm^k(x,\y)-\x\right)\right) d\theta \right)
\end{align*}
satisfies $|\partial_\x^\a\partial_\y^\b\partial_x^\g a(\x,\y,x)|\leq C_{\a,\b,\g}$. Finally we apply~\cite[Theorem~2.1]{As-F} to obtain the boundedness of the second term.

\eqref{lemm5.1.5} The first assertion follows from $s_k^\pm(x,\x)s_\ell^\pm(y,\x)=0$ for any $x$, $y$ and $\x$. For the second assertion, we set $\psi_k\in C^\infty(\T^d;M_n(\C))$ so that $\psi_k(\x)=P_k(\x)$ on $\supp \chi_k$. Then we use the equality $J_\pm^k=J_\pm^k \psi_k(D_x)$ and compactness of $[J_\pm^k,\psi_k(D_x)]$, which follows from~\eqref{lemm5.1.4}.
\end{proof}


Now we prove the existence of the following (inverse) local wave operators
\begin{align}
W^\pm(\JS)\, &{:=}\, \slim e^{itH} \JS e^{-itH_0} E_{H_0}(\Gamma), \label{local mWO}
\\
I^\pm(\JS)\, &{:=}\, \slim e^{itH_0} \JS^* e^{-itH} E_H^{\ac}(\Gamma), \label{i local mWO}
\end{align}
for $\JS=J_\#^k$ with $k=1,\dots,K$ and $\#\in\{+,-\}$. Note that, if $\JS$ is compact, then $W^\pm(\JS)=I^\pm(\JS)=0$.

We set $\tilde\chi_k \in C_c^\infty(\US_k)$ so that $\tilde\chi_k=1$ on $\supp\chi_k$. Since
\begin{align*}
(P_k \tilde\chi_k)(D_x) J_\#^k - J_\#^k = \left[(P_k \tilde\chi_k)(D_x), J_\#^k\right]
\end{align*}
is compact by Lemma~\ref{J lem}$\MK$(4), we have
\begin{align*}
W^\pm\left(J_\#^k\right) &= W^\pm\bigl(\left(P_k \tilde\chi_k\right)(D_x) J_\#^k\bigr), \\% = \slim e^{itH} (P_k \chi)(D_x) J_\#^k e^{-itH_0}.
I^\pm\left(J_\#^k\right) &= I^\pm\bigl(\left(P_k \tilde\chi_k\right)(D_x) J_\#^k\bigr),
\end{align*}
and thus it suffices to show the existence of~\eqref{local mWO} and~\eqref{i local mWO} for
\begin{align*}
\JS=(P_k \tilde\chi_k)(D_x) J_\#^k.
\end{align*}



\begin{lemm}\label{HJ-JH_0 lem}
\begin{multline*}
\bigl(H (P_k \tilde\chi_k)(D_x)J_\pm^k - (P_k \tilde\chi_k)(D_x)J_\pm^k H_0\bigr)u(x)\\
= (2\pi)^{-d}\int_{\T^d}\sum_{y\,\in\,\Z^d}e^{i\left(\f_\pm^k(x,\x)-y\cdot\xi\right)}a_\pm^k(x,\xi)u(y)d\xi,
\end{multline*}
where
\begin{multline}\label{s_a def}
a_\pm^k(x,\x)
= - i \y(x)\s_\pm'\bigl(\cos(x,\nabla_\x \l_k(\x))\bigr) \\
\frac{|\nabla_\x \l_k(\x)|^2 - |x|^{-2} (x\cdot\nabla_\x \l_k(\x))^2}{|x| |\nabla_\x \l_k(\x)|} P_k(\x) \chi_k(\x) 
 + r_\pm^k(x,\x) 
\end{multline}
and $|\partial_\x^\b r_\pm^k(x,\x)| \leq C_\b \la x \ra^{-\min(1+\rho,2)}$.
\end{lemm}

\begin{proofc}
\begin{step}\label{step1}
Let
\begin{align*}
g(x):=&\,(2\pi)^{-d}\int_{\T^d}e^{ix\cdot\x}H_0(\x)P_k(\x)\tilde\chi_k(\x)d\x \\
=&\,(2\pi)^{-d}\int_{\T^d}e^{ix\cdot\x}\lambda_k(\x)P_k(\x)\tilde\chi_k(\x)d\x.
\end{align*}
Then we learn
\begin{align*}
H_0 (P_k \tilde\chi_k)(D_x)J_\pm^k u(x) &= (2\pi)^{-d}\int_{\T^d}\sum_{y\,\in\,\Z^d}e^{i\left(\f_\pm^k(x,\x)-y\cdot\xi\right)}a_\pm^{k,1}(x,\xi)u(y)d\xi,
\end{align*}
where
\begin{align*}
a_\pm^{k,1}(x,\xi) &= \sum_{y\,\in\,\Z^d} g(y) e^{i\left(\f_\pm^k(x-y,\x)-\f_\pm^k(x,\x)\right)}s_\pm^k(x-y,\x) \\
&= \sum_{y\,\in\,\Z^d} g(y) e^{-iy\cdot\nabla_x\f_\pm^k(x,\x)} (1+R(x,y,\x)) s_\pm^k(x-y,\x),
\end{align*}
and
\begin{align*}
R(x,y,\x) := \exp \left[i\left(\f_\pm^k(x-y,\x)-\f_\pm^k(x,\x) + y\cdot\nabla_x\f_\pm^k(x,\x) \right)\right] -1.
\end{align*}
Since
\begin{multline*}
\left|\partial_\x^\b \left[\f_\pm^k(x-y,\x)-\f_\pm^k(x,\x) + y\cdot \nabla_x\f_\pm^k(x,\x) \right]\right| \\
\begin{aligned}
&= \left| y \cdot \int_0^1 \partial_\x^\b \left(\nabla_x\f_\pm^k(x,\x) - \nabla_x\f_\pm^k(x-\theta y, \x) \right) d\theta \right| \\
&= \left| y \cdot \int_0^1 \left(\int_0^1 \partial_\x^\b \nabla_x^2\f_\pm^k(x-\phi\theta y, \x) d\phi \right) \theta y d\theta \right| \\
&\leq C_\b \la x \ra^{-1-\rho} \jap{y}^{3+\rho},
\end{aligned}
\end{multline*}
we learn $|\partial_\x^\b R(x,y,\x)|\leq C_\b' \la x \ra^{-1-\rho} \jap{y}^{(3+\rho)\max\{1,|\b|\}}$, and thus
\begin{align*}
\left|\partial_\x^\b\sum_{y\,\in\,\Z^d} g(y) e^{-iy\cdot\nabla_x\f_\pm^k(x,\x)} R(x,y,\x) s_\pm^k(x-y,\x)\right|\leq C''_\b \la x \ra^{-1-\rho}.
\end{align*}
Furthermore, since~\eqref{s is in S^0} implies the similar inequality
\begin{multline*}
 \left|\partial_\x^\b \left[s_\pm^k(x-y,\x) - s_\pm^k(x,\x) + y \cdot \nabla_x s_\pm^k(x,\x)\right]\right|\\
= \left| y \cdot\int_0^1 \left(\int_0^1 \partial_\x^\b\nabla_x^2 s_\pm^k(x-\phi\theta y,\x) d\phi\right) \theta y d\theta \right|
\leq C_\b \jap{x}^{-2}\jap{y}^4,
\end{multline*}
we have
\begin{multline*}
\sum_{y\,\in\,\Z^d} g(y) e^{-iy\cdot\nabla_x\f_\pm^k(x,\x)}s_\pm^k(x-y,\x) \\
\begin{aligned}
&= \sum_{y\,\in\,\Z^d} g(y) e^{-iy\cdot\nabla_x\f_\pm^k(x,\x)} \left(s_\pm^k(x,\x) - y\cdot\nabla_x s_\pm^k(x,\x) \right) + O\left(\la x \ra^{-2}\right) \\
&= (\lambda_kP_k\tilde\chi_k)\left(\nabla_x\f_\pm^k(x,\x)\right) s_\pm^k(x,\x) - i \nabla_\x (\lambda_kP_k\tilde\chi_k)\left(\nabla_x\f_\pm^k(x,\x)\right)\cdot\nabla_x s_\pm^k(x,\x) \\
& \qquad + O\left(\la x \ra^{-2}\right).
\end{aligned}
\end{multline*}
Thus we obtain
\begin{multline*}
a_\pm^{k,1}(x,\xi) 
= (\lambda_kP_k\tilde\chi_k)\left(\nabla_x\f_\pm^k(x,\x)\right) s_\pm^k(x,\x) \\
 - i \nabla_\x (\lambda_kP_k\tilde\chi_k)\left(\nabla_x\f_\pm^k(x,\x)\right)\cdot\nabla_x s_\pm^k(x,\x)
 + O\left(\la x \ra^{-\min(1+\rho,2)}\right).
\end{multline*}

Similar computations imply that
\begin{align*}
V (P_k \tilde\chi_k)(D_x)J_\pm^k u(x) &= (2\pi)^{-d}\int_{\T^d}\sum_{y\,\in\,\Z^d}e^{i\left(\f_\pm^k(x,\x)-y\cdot\xi\right)}a_\pm^{k,2}(x,\xi)u(y)d\xi, \\
(P_k \tilde\chi_k)(D_x)J_\pm^k H_0 u(x) &= (2\pi)^{-d}\int_{\T^d}\sum_{y\,\in\,\Z^d}e^{i\left(\f_\pm^k(x,\x)-y\cdot\xi\right)}a_\pm^{k,3}(x,\xi)u(y)d\xi,
\end{align*}
where
\begin{align*}
&a_\pm^{k,2}(x,\xi) \\
=\,& V(x) \left((P_k\tilde\chi_k)\bigl(\nabla_x\f_\pm^k(x,\x)\bigr) s_\pm^k(x,\x) - i \nabla_\x (P_k\tilde\chi_k)\bigl(\nabla_x\f_\pm^k(x,\x)\bigr)\cdot\nabla_x s_\pm^k(x,\x) \right) \\
& \quad + O\left(\la x \ra^{-\rho-\min(1+\rho,2)}\right), \\
&a_\pm^{k,3}(x,\xi) \\
=\,& \lambda_k(\x) (P_k\tilde\chi_k)\bigl(\nabla_x\f_\pm^k(x,\x)\bigr) s_\pm^k(x,\x) - i \lambda_k(\x) \nabla_\x (P_k\tilde\chi_k)\bigl(\nabla_x\f_\pm^k(x,\x)\bigr)\cdot\nabla_x s_\pm^k(x,\x) \\
& \quad + O\left(\la x \ra^{-\min(1+\rho,2)}\right).
\end{align*}
\end{step}
\begin{step}\label{step2}
Step~\ref{step1} implies
\begin{multline*}
a_\pm^k(x,\x)\\
\begin{aligned}
=&\, (P_k\tilde\chi_k)\bigl(\nabla_x\f_\pm^k(x,\x)\bigr) s_\pm^k(x,\x) \Bigl(\lambda_k\bigl(\nabla_x\f_\pm^k(x,\x)\bigr)+V(x)-\lambda_k(\x)\Bigr) \\
&-i \nabla_\x (P_k\tilde\chi_k)\bigl(\nabla_x\f_\pm^k(x,\x)\bigr)\cdot \nabla_x s_\pm^k(x,\x) \Bigl(\lambda_k\bigl(\nabla_x\f_\pm^k(x,\x)\bigr)+V(x)-\lambda_k(\x)\Bigr) \\
&-i (P_k\tilde\chi_k)\bigl(\nabla_x\f_\pm^k(x,\x)\bigr)\nabla_\x \lambda_k\bigl(\nabla_x\f_\pm^k(x,\x)\bigr)\cdot\nabla_x s_\pm^k(x,\x) \\
& +O\left(\la x \ra^{-\min(1+\rho,2)}\right).
\end{aligned}
\end{multline*}
The first and second terms are of order $\jap{x}^{-1-\rho}$ by~\eqref{eikonal} and~\eqref{short range}. Moreover simple computations imply that, setting $v:=\nabla_\x\lambda_k(\x)$,
\begin{align*}
\nabla_x s_\pm^k(x,\x) 
= \eta(x)\s_\pm'\bigl(\cos(x,v) \bigr) \left(\frac{1}{|x||v|}v - \frac{x\cdot v}{|x|^3|v|}x \right) P_k(\x) \chi_k(\x) +O\left(\jap{x}^{-\infty}\right),
\end{align*}
and therefore
\begin{align*}
a_\pm^k(x,\x) 
=& -i (P_k\tilde\chi_k)(\x)\nabla_\x \lambda_k(\x)\cdot\nabla_x s_\pm^k(x,\x) + O\left(\la x \ra^{-\min(1+\rho,2)}\right) \\
=& -i \eta(x)\s_\pm'\bigl(\cos(x,v) \bigr) \left(\frac{|v|}{|x|}-\frac{(x\cdot v)^2}{|x|^3|v|}\right) P_k(\x) \chi_k(\x) + O\left(\la x \ra^{-\min(1+\rho,2)}\right).
\end{align*}
Here we have used~\eqref{phase estimate} in the first equality to replace $\nabla_x\f_\pm^k(x,\x)$ by $\x$.\qedhere
\end{step}\let\qed\relax
\end{proofc}



\begin{prop}\label{existence of local WO}
For any $k=1,\dots,K$, there exist the limits~\eqref{local mWO} and~\eqref{i local mWO} with $\JS=J_\pm^k$.
\end{prop}

\begin{proof}
We only prove the existence of~\eqref{local mWO}, since the other is done in the same way.

We may assume $\rho<1$ without loss of generality. The standard argument of existence of (modified) wave operators (see, e.g., \cite[Lemmas~10.2.1 and 10.2.2, Theorem~0.5.4]{Y} and~\cite[Theorem~XIII. 24]{R-S4}) implies that it suffices to prove that $H (P_k \tilde\chi_k)(D_x)J_\pm^k - (P_k \tilde\chi_k)(D_x)J_\pm^k H_0$ is a finite sum of the form $G_j^* B_j G'_j$ with $G_j$ (resp.\ $G'_j$) being $H$-(resp.\ $H_0$-) smooth in $\Gamma$ and $B_j\in\BS(\HS)$.

We set
\begin{align*}
a_j^k(x,\x) &= \y(x)|x|^{-\frac{1}{2}}\left(\partial_{\x_j}\l_k(\x)-|x|^{-2} x_j \left(x\cdot\nabla_\x\l_k(\x)\right)\right) P_k(\x) \tilde\chi_k(\x), \\
b_{\pm}^k(x,\x) &= -i \y(x)\s_\pm'\bigl(\cos(x,\nabla_\x\lambda_k(\x)) \bigr) P_k(\x) \chi_k(\x).
\end{align*}
Then we observe that
\begin{align*}
a_j^k(x,D_x)=\y(x)|x|^{-\frac12}\nabla_{k,j}^\perp,
\end{align*}
where $\nabla_{k,j}^\perp$ is as in Proposition~\ref{prop_rad_est}. Moreover we have by the definition~\eqref{s_a def} of $a_\pm^k(x,\x)$
\begin{align*}
a_\pm^k(x,\x)=b_\pm^k(x,\x) \sum_{j=1}^d a_j^k(x,\x)^2 + r_\pm^k(x,\x),
\end{align*}
where $\partial_\x^\a r_\pm^k(x,\x) = O(\jap{x}^{-2})$.

We take functions $\tilde{\tilde{\chi}}_k \in C_c^\infty(\US_k)$ and $\tilde\s_\pm(\theta) \in C^\infty(\R)$ such that
\begin{align*}
&\tilde\s_\pm(\theta)=
\begin{alcases}
1 & \quad \text{if } \pm\theta \geq -\frac{1}{2}, \\ 0 & \quad \text{if } \pm\theta \geq -\frac{3}{4},
\end{alcases}
\\
&\tilde{\tilde\chi}_k(\x) = 1, \quad \x\in \supp\tilde\chi_k.
\end{align*}
We set
\begin{align*}
\tilde s^k(x,\x) &= \eta(x) P_k(\x) \tilde{\tilde\chi}_k(\x), \\
\tilde \f_{\pm}^k(x,\x) &= \y(x)\tilde\s_\pm\bigl(\cos(x,\nabla\lambda_k(\x))\bigr) \f_{\pm}^k(x,\x) + \Bigl(1 - \y(x)\tilde\s_\pm\bigl(\cos(x,\nabla\lambda_k(\x))\bigr) \Bigr) x\cdot\x,
\end{align*}
and
\begin{align*}
\tilde J_\pm^k u(x) &= (2\pi)^{-\frac{d}{2}} \int_{\T^d} e^{i\tilde\f_{\pm}^k(x,\x)} \tilde s^k(x,\x) \FS u(\x) d\x, \\
A_{\pm,j}^ku(x)&=(2\pi)^{-\frac{d}{2}}\int_{\T^d} e^{i\tilde\f_\pm^k(x,\x)} a_j^k(x,\x) \FS u(\x)d\x, \\
C_{\pm,j}^ku(x) &= (2\pi)^{-\frac{d}{2}}\int_{\T^d} e^{i\f_\pm^k(x,\x)} b_\pm^k(x,\x)a_j^k(x,\x)^2 \FS u(\x)d\x.
\end{align*}
Then it follows from the same argument as Lemma~\ref{J lem}$\MK$\eqref{lemm5.1.2} that
\begin{align*}
&\tilde J_\pm^k (\tilde J_\pm^k)^* = \tilde s^k(x,D_x)^2 +R_{\pm,j,1}^k, \\
&(\tilde J_\pm^k)^*A_{\pm,j}^k = a_j^k(x,D_x)+R_{\pm,j,2}^k, \\
&(\tilde J_\pm^k)^*C_{\pm,j}^k = a_j^k(x,D_x) b_\pm^k(x,D_x) a_j^k(x,D_x)+R_{\pm,j,3}^k,
\end{align*}
where $\jap{x}^{\frac{1+\rho}{2}} R_{\pm,j,\ell}^k \jap{x}^{\frac{1+\rho}{2}} \in \BS(\HS)$, $\ell=1$, $2$, $3$. Moreover we learn by the argument in Lemma~\ref{J lem}$\MK$(4) that
\begin{align*}
&\tilde s^k(x,D_x)^2 A_{\pm,j}^k = A_{\pm,j}^k + R_{\pm,j,4}^k, \\
&\tilde s^k(x,D_x)^2 C_{\pm,j}^k = C_{\pm,j}^k + R_{\pm,j,5}^k,
\end{align*}
where $\jap{x}^{\frac{1+\rho}{2}} R_{\pm,j,\ell}^k \jap{x}^{\frac{1+\rho}{2}}\in \BS(\HS)$, $\ell=4$, $5$. Thus we have, modulo operators of the form $\jap{x}^{-\frac{1+\rho}{2}}B\jap{x}^{-\frac{1+\rho}{2}}$ with $B\in\BS(\HS)$,
\begin{align*}
H (P_k \tilde\chi_k)(D_x)J_\pm^k - (P_k \tilde\chi_k)(D_x)J_\pm^k H_0 
&\equiv \sum_{j=1}^d C_{\pm,j}^k \\
&\equiv \sum_{j=1}^d \tilde s^k(x,D_x)^2 C_{\pm,j}^k \\
&\equiv \sum_{j=1}^d \tilde J_\pm^k (\tilde J_\pm^k)^* C_{\pm,j}^k \\
&\equiv \sum_{j=1}^d \tilde J_\pm^k a_j^k(x,D_x) b_\pm^k(x,D_x) a_j^k(x,D_x) \\
&\equiv \sum_{j=1}^d \tilde J_\pm^k (\tilde J_\pm^k)^* A_{\pm,j}^k b_\pm^k(x,D_x) a_j^k(x,D_x) \\
&\equiv \sum_{j=1}^d \tilde s^k(x,D_x)^2 A_{\pm,j}^k b_\pm^k(x,D_x) a_j^k(x,D_x) \\
&\equiv \sum_{j=1}^d A_{\pm,j}^k b_\pm^k(x,D_x) a_j^k(x,D_x).
\end{align*}

Since $b_\pm^k(x,D_x) \in \BS(\HS)$ and Proposition~\ref{prop_rad_est} implies $a_j^k(x,D_x)$ is $H_0$-smooth on $\Gamma$, it remains to prove that $A_{\pm,j}^k$ is $H$-smooth on $\Gamma$. However, the proof is completed if we observe that $a_j^k(x,D_x)$ and $\jap{x}^{\frac{1+\rho}{2}}$ are $H$-smooth on $\Gamma$ and that
\begin{align*}
\left(A_{\pm,j}^k\right)^* A_{\pm,j}^k=a_j^k(x,D_x)^*a_j^k(x,D_x) + R''_j,
\end{align*}
where $\jap{x}^{\frac{1+\rho}{2}}R_j''\jap{x}^{\frac{1+\rho}{2}}\in\BS(\HS)$.
\end{proof}



\section{Proof of Theorem~\ref{main}} \label{proof section}


We set
\begin{align}\label{time-independent modifiers}
J_\pm := \sum_{k=1}^K J_\pm^k,
\end{align}
where $J_\pm^k$'s are given by~\eqref{IK1}. Then Proposition~\ref{existence of local WO} implies the existence of the modified wave operators~\eqref{modified WOs}. The proof of the intertwining property is skipped since it is easily proved.



\begin{prop}\label{nonstationary prop}
$W^\pm(J_\mp) = I^\pm(J_\mp) = 0$.
\end{prop}

\begin{proof}
For the first assertion, it suffices to prove $\lim_{t\to\pm\infty} J_\mp^k e^{-itH_0}u=0$ for any $u$ satisfying
\begin{align*}
(P_k\chi_k)(D_x)u=u.
\end{align*}
We easily see that
\begin{align*}
J_\mp^k e^{-itH_0} u(x) 
=(2\pi)^{-\frac{d}{2}} \int_{\T^d} e^{i\left(\f_\mp^k(x,\x)-t\lambda_k(\x)\right)} \eta(x)\s_\mp\bigl(\cos(x,\nabla\lambda_k(\x))\bigr) \FS u(\x)d\x.
\end{align*}
The estimate~\eqref{phase estimate} and the conditions~\eqref{eta 01} and~\eqref{sigma 01} imply there is a constant $c>0$ such that on the support of the integrand
\begin{align*}
\left|\nabla_\x \f_\mp^k(x,\x)-t\nabla\lambda_k(\x)\right|
&\geq |x-t\nabla\lambda_k(\x)| - \left|x-\nabla_\x\f_\mp^k(x,\x)\right| \\
&\geq \sqrt{\frac{1-\cos(x,\pm\nabla\lambda_k(\x))}{2}}|x| |t\nabla\lambda_k(\x)| - C\jap{x}^{1-\rho} \\
&\geq c \big(|x| + |t| |\nabla\lambda_k(\x)|\big)
\end{align*}
for sufficiently large $\pm t \geq 0$. The non-stationary phase method implies that
\begin{align*}
\left|J_\mp^k e^{-itH_0} u(x)\right|\leq C_{N}(1+|x|+|t|)^{-N}, \quad x\in\Z^d,\ \pm t \geq0,
\end{align*}
for any $N\geq1$. Thus we obtain $\|W^\pm(J_\mp)u\| = 0$.

For the other assertion $I^\pm(J_\mp) = 0$, the intertwining property implies
\begin{align*}
I^\pm(\JS)=I^\pm(\JS) E_H(\Gamma) = E_{H_0}(\Gamma) I^\pm(\JS).
\end{align*}
Thus we learn that for any $v\in\HS$
\begin{align*}
\left(I^\pm(J_\mp)u, v\right) &= \left(E_{H_0}(\Gamma) I^\pm(J_\mp)u, v\right) \\
&= \lim_{t\to\pm\infty} \left(e^{itH_0} J_\mp^* e^{-itH} E_{H}^{\ac}(\Gamma) u, E_{H_0}(\Gamma) v\right) \\
&= \lim_{t\to\pm\infty} \left(E_{H}^{\ac}(\Gamma) u, e^{itH} J_\mp e^{-itH_0} E_{H_0}(\Gamma) v\right) \\
&= \left(E_{H}^{\ac}(\Gamma) u, W^\pm(J_\mp) v\right) \\
&= 0
\end{align*}
by the first assertion.
\end{proof}

\begin{prop}
For any $u\in\HS$,
\begin{align}
\left\|W^\pm(J_\pm)u\right\| =& \left\|E_{H_0}(\Gamma)u\right\|, \label{p isom}
\\
\left\|I^\pm(J_\pm)u\right\| =& \left\|E_{H}^{\ac}(\Gamma)u\right\|. \label{i p isom}
\end{align}
\end{prop}

\begin{proof}
We learn
\begin{align*}
\left\|W^\pm(\JS)u\right\|^2 = \lim_{t\to\pm\infty} \left\| \JS e^{-itH_0}E_{H_0}(\Gamma)u \right\|^2 = \lim_{t\to\pm\infty} \left(u_t, \JS^*\JS u_t \right),
\end{align*}
where $u_t:=e^{-itH_0}E_{H_0}(\Gamma)u$. Thus Lemmas~\ref{adjoint lemma}, \ref{composition lemma}, \ref{J lem}$\MK$\eqref{lemm5.1.2}, \eqref{lemm5.1.5} and~\eqref{Rep of E_H0}, \eqref{chi 1}, \eqref{sigma sum1} imply
\begin{multline*}
\left\|W^\pm(J_+)u\right\|^2 + \left\|W^\pm(J_-)u\right\|^2 \\
\begin{aligned}
&= \lim_{t\to\pm\infty} \left(u_t, \left(J_+^*J_+ + J_-^*J_-\right) u_t \right) \\
&= \lim_{t\to\pm\infty} \left(u_t, \left(\;\sum_{k=1}^K \left(J_+^k\right)^* J_+^k + \left(J_-^k\right)^* J_-^k\right) u_t \right) \\
&= \lim_{t\to\pm\infty} \left(u_t,
\left(\;\sum_{k=1}^K s_+^k(x,D_x) s_+^k(x,D_x)^* + s_-^k(x,D_x) s_-^k(x,D_x)^* \right) u_t \right) \\
&= \lim_{t\to\pm\infty} \left(u_t,
\y(x)^2 \sum_{k=1}^K\left(P_k\chi_k^2\right)(D_x) u_t \right) \\
&= \lim_{t\to\pm\infty} \left(u_t,
\y(x)^2 u_t \right) \\
&= \left\|E_{H_0}(\Gamma)u \right\|^2.
\end{aligned}
\end{multline*}
Here we have used~\eqref{Rep of E_H0} and~\eqref{eta 01} to obtain $\sum_{k=1}^K (P_k\chi_k^2)(D_x) E_{H_0}(\Gamma) = E_{H_0}(\Gamma)$ and compactness of $1-\y(x)^2$. Therefore we have the first equality~\eqref{p isom} by Proposition~\ref{nonstationary prop}.

The other equality~\eqref{i p isom} is obtained by the similar argument and the compactness of $\psi(H)-\psi(H_0)$ for $\psi\in C_c^\infty(\R)$.
\end{proof}



It remains to prove the completeness of~\eqref{modified WOs}. However it is proved by the existence of $I^\pm(J_\pm)$ and~\eqref{i p isom}.



\bibliography{tadano}
\end{document}
